%%% FIELD ass-strong-condition
\(P^{\star}(D_1\ge D_0)=1\).
%%% FIELD ass-conditional-randomization
\(Z\perp_{P^{\dagger}}(D^{\mathrm{pot}},Y^{\mathrm{pot}})\mid X\).
%%% FIELD ass-conditional-overlap
\(0<P^{\dagger}(Z=1\mid X=x)<1\) for every \(x\in\mathcal X\) with \(P^{\dagger}(X=x)>0\).
%%% FIELD ass-linear-response-restriction
For every \(x\in\mathcal X\) with \(P(X=x)>0\), the response-mass vector \(\bigl(P^{\dagger}(D_0=i,D_1=j\mid X=x)\bigr)_{0\le i,j\le L}\) belongs to the researcher-chosen nonempty polytope \(\mathcal R_{T,x}\).
%%% FIELD ass-general-stayer-floor
\[\inf_{\gamma\in\Gamma_{\mathcal R_T}(\mu_x)}\mathsf S(\gamma)\ge\kappa\qquad\text{for every }x\in\mathcal X.\]
%%% FIELD def-model-class-replacement
\(\mathcal M_L=\{P^{\star}: P^{\star}\text{ satisfies }\mathrm{ass{:}random\mbox{-}assignment}\text{ and }\mathrm{ass{:}strong\mbox{-}tier\mbox{-}monotonicity}\text{ for the fixed }L\}\).
%%% FIELD def-linear-response-model
\[\mathfrak M_{L,\mathcal R_T}=\left\{P^{\dagger}:P^{\dagger}\text{ is a law of }(X,Z,D^{\mathrm{pot}},Y^{\mathrm{pot}})\text{ satisfying }\mathrm{ass{:}conditional\mbox{-}random\mbox{-}assignment},\ \mathrm{ass{:}conditional\mbox{-}assignment\mbox{-}overlap},\text{ and }\mathrm{ass{:}linear\mbox{-}response\mbox{-}type\mbox{-}restriction}\right\}.\]
%%% FIELD def-conditional-targets
Let \(\mathcal X\) be a finite nonempty support and let \(\omega_x=P^{\dagger}(X=x)>0\), with \(\sum_{x\in\mathcal X}\omega_x=1\).  For \(z\in\{0,1\}\), define the conditional arm-state probabilities
\[\mu_{zx}(0^{\star})=P^{\dagger}(D=0\mid Z=z,X=x),\qquad
  \mu_{zx}(d,y)=P^{\dagger}(D=d,Y^{\mathrm{obs}}=y\mid Z=z,X=x)\]
for \(d=1,\ldots,L\) and \(y\in\{0,1\}\), and put \(\mu_x=(\mu_{0x},\mu_{1x})\).  Whenever the conditioning event in the denominator has positive probability, define
\[H_x^{\mathcal R}=P^{\dagger}(Y_{0,D_0}=1,Y_{1,D_0}=0\mid D_0=D_1>0,X=x).\]
For a compatible conditional arm-cell array \(\mu=(\mu_x)_{x\in\mathcal X}\), let \(\mathfrak P_{\mathcal R_T}(\mu,\omega)\) be the fiber of laws \(Q\in\mathfrak M_{L,\mathcal R_T}\) that have \(Q(X=x)=\omega_x\) and induce \(\mu\) under the observation map.  The fixed-stratum identified set is
\[\mathcal I_{\mathcal R_T,x}(\mu_x)=\{H_x^{\mathcal R}(Q):Q\in\mathfrak P_{\mathcal R_T}(\mu,\omega)\}.\]
Writing \(S_x(Q)=Q(D_0=D_1>0\mid X=x)\) and \(N_x(Q)=Q(D_0=D_1>0,Y_{0,D_0}=1,Y_{1,D_0}=0\mid X=x)\), define the finite-strata target and its identified set by
\[H_{\mathcal X}^{\mathcal R}(Q)=\frac{\sum_{x\in\mathcal X}\omega_xN_x(Q)}{\sum_{x\in\mathcal X}\omega_xS_x(Q)},\qquad
  \mathcal I_{\mathcal R_T,\mathcal X}(\mu,\omega)=\{H_{\mathcal X}^{\mathcal R}(Q):Q\in\mathfrak P_{\mathcal R_T}(\mu,\omega)\},\]
whenever the displayed denominator is positive.
%%% FIELD def-linear-response-flow
Let
\[\mathcal E=\{0^{\star}\}\cup\{(d,y):d=1,\ldots,L,\ y\in\{0,1\}\},\qquad
  \lambda(0^{\star})=0,\quad \lambda(d,y)=d.\]
Let \(m_T=(L+1)^2\).  For each \(x\in\mathcal X\), fix finite row counts \(k_x^=,k_x^{\le}\), matrices \(A_x\in\mathbb R^{k_x^=\times m_T}\) and \(B_x\in\mathbb R^{k_x^{\le}\times m_T}\), and vectors \(a_x\in\mathbb R^{k_x^=}\) and \(b_x\in\mathbb R^{k_x^{\le}}\).  A linear response-type restriction is a nonempty polytope
\[\mathcal R_{T,x}=\{p\in\mathbb R_+^{(L+1)^2}:\mathbf 1^{\mathsf T}p=1,\ A_xp=a_x,\ B_xp\le b_x\}.\]
Support restrictions are included by rows such as \(p_{ij}=0\).  An expanded-state arc \((e_0,e_1)\) is allowed when the layer pair \((\lambda(e_0),\lambda(e_1))\) is not forced to have zero response mass by \(\mathcal R_{T,x}\).  For a nonnegative array \(\gamma=(\gamma(e_0,e_1))\) on the allowed arcs, set
\[p_{ij}(\gamma)=\sum_{e_0:\lambda(e_0)=i}\sum_{e_1:\lambda(e_1)=j}\gamma(e_0,e_1).\]
The conditional expanded-flow polytope is
\[\Gamma_{\mathcal R_T}(\mu_x)=\left\{\gamma\ge0:
\sum_{e_1}\gamma(e_0,e_1)=\mu_{0x}(e_0),\ 
\sum_{e_0}\gamma(e_0,e_1)=\mu_{1x}(e_1),\ 
p(\gamma)\in\mathcal R_{T,x}\right\}.\]
Define the stayer and harm arc costs
\[\sigma(e_0,e_1)=\mathbf 1\{\lambda(e_0)=\lambda(e_1)>0\},\qquad
  \eta(e_0,e_1)=\sum_{d=1}^{L}\mathbf 1\{e_0=(d,1),\ e_1=(d,0)\},\]
and the linear functionals
\[\mathsf S(\gamma)=\sum_{e_0,e_1}\sigma(e_0,e_1)\gamma(e_0,e_1),\qquad
  \mathsf N(\gamma)=\sum_{e_0,e_1}\eta(e_0,e_1)\gamma(e_0,e_1).\]
Let \(J_x\) be the number of allowed arcs.  After flattening those arcs and collecting the finitely many equality and inequality constraints, there are finite row counts \(k_x^G,k_x^F\), matrices \(G_x\in\mathbb R^{k_x^G\times J_x}\) and \(F_x\in\mathbb R^{k_x^F\times J_x}\), and vectors \(g_x(\mu_x)\in\mathbb R^{k_x^G}\) and \(f_x\in\mathbb R^{k_x^F}\) such that
\[\Gamma_{\mathcal R_T}(\mu_x)=\{\gamma\ge0:G_x\gamma=g_x(\mu_x),\ F_x\gamma\le f_x\}.\]
The margin equations in \(G_x\gamma=g_x(\mu_x)\) include \(\mathbf 1^{\mathsf T}\gamma=1\).  Under Lee's binary positive-selection monotonicity alone, \(p_{i0}=0\) for \(i>0\), while both directions \((i,j)\) and \((j,i)\) remain allowed for distinct positive \(i,j\).  Under KMV Strong Monotonicity,
\[\mathcal R_T^{\uparrow}=\{p:p_{ij}=0\text{ whenever }j<i\},\]
so the allowed layer arcs are exactly upper triangular.
%%% FIELD lem-kmv-framework-statement
Kroft, Mourifi\'e, and Vayalinkal formulate researcher-chosen linear response-type restrictions \(\mathcal R_T\), explicitly include \(D_1\ge D_0\) as their Strong Monotonicity restriction, and state their main sharp-bounds theorem for the response-type framework with the chosen \(\mathcal R_T\).
%%% FIELD lem-general-decoder-statement
Fix \(x\in\mathcal X\).  Under conditional random assignment, conditional assignment overlap, and the linear response-type restriction \(\mathcal R_{T,x}\), mapping a full-data law to the conditional joint law \(\gamma\) of its two realized arm states maps exactly onto \(\Gamma_{\mathcal R_T}(\mu_x)\).  Conversely, every \(\gamma\in\Gamma_{\mathcal R_T}(\mu_x)\) has a full-data completion that induces \(\mu_x\), satisfies the response-type restriction, and obeys \(\mathsf S(\gamma)=P^{\dagger}(D_0=D_1>0\mid X=x)\) and \(\mathsf N(\gamma)=P^{\dagger}(D_0=D_1>0,Y_{0,D_0}=1,Y_{1,D_0}=0\mid X=x)\).  Thus the correspondence is a bijection after full-data laws are quotiented by their joint realized-arm-state distribution, including when some observable cells are zero.
%%% FIELD lem-general-decoder-proof
Fix \(x\) with \(\omega_x>0\).  For \(z\in\{0,1\}\), define the potential realized arm state
\[E_z=\begin{cases}
0^{\star},&D_z=0,\\
(D_z,Y_{z,D_z}),&D_z>0.
\end{cases}\]
This definition uses only potential outcomes that are observed under the corresponding assignment.  Given an admissible full-data law, put
\[\gamma(e_0,e_1)=P^{\dagger}(E_0=e_0,E_1=e_1\mid X=x).\]
Conditional random assignment implies
\[P^{\dagger}(E_z=e\mid X=x)=P^{\dagger}(E_z=e\mid Z=z,X=x)=\mu_{zx}(e),\]
where the second conditional probability is well defined by conditional overlap and the last equality follows from the observation map.  Hence the row and column margins of \(\gamma\) are \(\mu_{0x}\) and \(\mu_{1x}\).  Moreover,
\[p_{ij}(\gamma)=P^{\dagger}(D_0=i,D_1=j\mid X=x),\]
so the maintained response-type restriction gives \(p(\gamma)\in\mathcal R_{T,x}\).  Therefore \(\gamma\in\Gamma_{\mathcal R_T}(\mu_x)\).

For the reverse map, fix \(\gamma\in\Gamma_{\mathcal R_T}(\mu_x)\) and draw \((E_0,E_1)\) with this mass function conditional on \(X=x\).  Set \(D_z=\lambda(E_z)\).  If \(E_z=(d,y)\), set \(Y_{z,d}=y\); set every other, unused coordinate of \(Y^{\mathrm{pot}}\) to zero.  This produces a well-defined array in \(\{0,1\}^{2L}\), also when \(E_z=0^{\star}\).  Its conditional response masses are \(p(\gamma)\), hence belong to \(\mathcal R_{T,x}\), and its two potential realized-state margins are exactly the two margins of \(\gamma\).  Combining these conditional constructions over the finitely many strata with masses \(\omega_x\), and then drawing \(Z\) conditionally independently with its observed propensity in \((0,1)\), gives a normalized law \(P^{\dagger}\) satisfying conditional random assignment and overlap and inducing the prescribed observed arm cells.

Finally, by the definitions of \(\sigma\) and \(\eta\),
\[\mathsf S(\gamma)=\sum_{d=1}^{L}\sum_{y_0,y_1\in\{0,1\}}P^{\dagger}(E_0=(d,y_0),E_1=(d,y_1)\mid X=x)
=P^{\dagger}(D_0=D_1>0\mid X=x),\]
and
\[\mathsf N(\gamma)=\sum_{d=1}^{L}P^{\dagger}(E_0=(d,1),E_1=(d,0)\mid X=x),\]
which is the asserted conditional harm mass.  The construction is a right inverse for the realized-state joint distribution.  Unused potential outcomes are not identified, so quotienting by them is necessary and sufficient for bijectivity.  No division was used, and zero cells simply force the corresponding marginal and incident flow to be zero.
%%% FIELD thm-general-program-statement
Let \(\mathcal X\) be finite and let \(\mu=(\mu_x)_{x\in\mathcal X}\) be compatible in the sense that \(\Gamma_{\mathcal R_T}(\mu_x)\ne\varnothing\) for every \(x\).  Under the conditional random-assignment, conditional-overlap, linear-response-type, and positive-stayer-floor conditions, the sharp fixed-stratum identified set is
\[\mathcal I_{\mathcal R_T,x}(\mu_x)=\left[
\min_{\gamma\in\Gamma_{\mathcal R_T}(\mu_x)}\frac{\mathsf N(\gamma)}{\mathsf S(\gamma)},
\max_{\gamma\in\Gamma_{\mathcal R_T}(\mu_x)}\frac{\mathsf N(\gamma)}{\mathsf S(\gamma)}
\right].\]
Write a flattened version of the flow polytope as \(\{\gamma\ge0:G_x\gamma=g_x,\ F_x\gamma\le f_x\}\), and write \(\sigma_x\) and \(\eta_x\) for its stayer and harm cost vectors.  The two endpoints are equivalently the minimum and maximum of \(\eta_x^{\mathsf T}\xi\), respectively, over the same explicit Charnes--Cooper feasible set
\[G_x\xi=tg_x,\qquad F_x\xi\le tf_x,\qquad \sigma_x^{\mathsf T}\xi=1,\qquad \xi\ge0,\quad 0\le t\le\kappa^{-1}.\]
Both extrema and every value between them are attained by full-data laws.

For the finite-strata target, set \(\boldsymbol\Gamma=\prod_{x\in\mathcal X}\Gamma_{\mathcal R_T}(\mu_x)\).  Its sharp identified set is
\[\mathcal I_{\mathcal R_T,\mathcal X}(\mu,\omega)=\left[
\min_{(\gamma_x)\in\boldsymbol\Gamma}\frac{\sum_x\omega_x\mathsf N(\gamma_x)}{\sum_x\omega_x\mathsf S(\gamma_x)},
\max_{(\gamma_x)\in\boldsymbol\Gamma}\frac{\sum_x\omega_x\mathsf N(\gamma_x)}{\sum_x\omega_x\mathsf S(\gamma_x)}
\right].\]
These endpoints are the minimum and maximum of \(\sum_x\omega_x\eta_x^{\mathsf T}\xi_x\) over variables \(0\le t\le\kappa^{-1}\) and \(\xi_x\ge0\) satisfying
\[G_x\xi_x=tg_x,\qquad F_x\xi_x\le tf_x\quad(x\in\mathcal X),\qquad
  \sum_{x\in\mathcal X}\omega_x\sigma_x^{\mathsf T}\xi_x=1.\]
Thus the broader response-type characterization is computable by two finite linear programs; no ordered-support or closed-form claim is made for a general \(\mathcal R_T\).
%%% FIELD thm-general-program-proof
We distinguish the causal targets from the observable optimization maps.  At a fixed stratum \(x\), the decoder lemma gives both inclusions
\[\{(S_x(Q),N_x(Q)):Q\in\mathfrak P_{\mathcal R_T}(\mu,\omega)\}
=\{(\mathsf S(\gamma),\mathsf N(\gamma)):\gamma\in\Gamma_{\mathcal R_T}(\mu_x)\}.\]
Indeed, the forward inclusion sends each full-data law to its realized-state coupling, and the reverse inclusion is the explicit full-data completion.  The fixed-stratum version of the positive-stayer-floor condition gives \(\mathsf S(\gamma)\ge\kappa>0\) for every feasible \(\gamma\).  Consequently division is valid, and the definition of \(H_x^{\mathcal R}\) yields, rather than assumes,
\[\mathcal I_{\mathcal R_T,x}(\mu_x)
=\left\{\frac{\mathsf N(\gamma)}{\mathsf S(\gamma)}:\gamma\in\Gamma_{\mathcal R_T}(\mu_x)\right\}.\]

The set \(\Gamma_{\mathcal R_T}(\mu_x)\) is the intersection of finitely many closed affine half-spaces with the nonnegative unit-mass simplex.  Compatibility makes it nonempty; hence it is compact and convex.  Since \(\mathsf S\ge\kappa\), the ratio \(\mathsf N/\mathsf S\) is continuous on this polytope.  It therefore attains a minimum and maximum.  Convexity makes the polytope connected, and the continuous image of a connected set in \(\mathbb R\) is an interval.  Thus every intermediate ratio is also attained.  Applying the decoder again converts every optimizing or intermediate flow into an agreeing full-data law, which proves sharpness rather than merely an outer bound.

For completeness, we verify the advertised linearization equation by equation.  Given a feasible \(\gamma\), define
\[t=\frac{1}{\sigma_x^{\mathsf T}\gamma},\qquad \xi=t\gamma.\]
Positivity of the denominator gives \(t>0\).  Multiplication of the equality and inequality systems by \(t\) gives
\[G_x\xi=tg_x,\qquad F_x\xi\le tf_x,\qquad \xi\ge0,\qquad
  \sigma_x^{\mathsf T}\xi=1,\]
and the objective satisfies
\[\eta_x^{\mathsf T}\xi
=\frac{\eta_x^{\mathsf T}\gamma}{\sigma_x^{\mathsf T}\gamma}
=\frac{\mathsf N(\gamma)}{\mathsf S(\gamma)}.\]
The stayer floor also gives \(t\le\kappa^{-1}\).
Conversely, the scaled margin equations include \(\mathbf 1^{\mathsf T}\xi=t\).  Since \(0\le\sigma_x\le\mathbf 1\) coordinatewise and \(\sigma_x^{\mathsf T}\xi=1\),
\[t=\mathbf 1^{\mathsf T}\xi\ge\sigma_x^{\mathsf T}\xi=1,\]
so \(t>0\).  Setting \(\gamma=\xi/t\) then recovers every original constraint, gives \(\sigma_x^{\mathsf T}\gamma=1/t\), and preserves the displayed objective equality.  This is a bijection of feasible points modulo the deterministic scaling, proving equivalence of both the minimization and maximization programs.  The homogenizing substitution is the Charnes--Cooper technique (Charnes and Cooper, 1962, pp.~181--186, DOI \(10.1002/\mathrm{nav}.3800090303\)); the specialized equivalence needed here has just been derived directly.

For finitely many strata, apply the decoder independently conditional on each \(x\) and combine the resulting conditional laws with masses \(\omega_x\).  This gives
\[\mathcal I_{\mathcal R_T,\mathcal X}(\mu,\omega)
=\left\{\frac{\sum_x\omega_x\mathsf N(\gamma_x)}{\sum_x\omega_x\mathsf S(\gamma_x)}:
(\gamma_x)_x\in\boldsymbol\Gamma\right\}.\]
The assumed floor makes the weighted denominator at least \(\kappa\).  The finite product \(\boldsymbol\Gamma\) is nonempty, compact, and convex, so the same continuity and connectedness argument proves attainment and interval filling.  Finally set \(\xi_x=t\gamma_x\), where
\[t=\left(\sum_x\omega_x\sigma_x^{\mathsf T}\gamma_x\right)^{-1}.\]
The aggregate stayer floor gives \(t\le\kappa^{-1}\), so this gives exactly the displayed stacked constraints and weighted objective.  Conversely, each scaled conditional margin system gives \(\mathbf 1^{\mathsf T}\xi_x=t\); because \(\sum_x\omega_x=1\), the weighted normalization implies \(t\ge1\).  Division by \(t\) therefore recovers one feasible \(\gamma_x\) in every stratum.  This proves the aggregate linear programs and completes the characterization.
%%% FIELD prop-ordered-collapse-statement
Suppose \(\mathcal R_{T,x}=\mathcal R_T^{\uparrow}\) for every \(x\in\mathcal X\).  Decode from \(\mu_x\) the conditional primitives \(\ell_{xd},u_{xd},a_{xd},b_{xd}\) by applying \(\mathrm{def{:}observable\mbox{-}primitives}\) within stratum \(x\), and suppose every conditional observable law belongs to \(\mathcal D_{L,\kappa}\).  Then the broader fixed-stratum flow endpoints equal
\[H_{-,x}=\frac{\sum_{d=1}^{L}(\ell_{xd}-a_{xd})_+}{\sum_{d=1}^{L}\min\{u_{xd},\max\{\ell_{xd},a_{xd}\}\}},\qquad
H_{+,x}=\frac{\sum_{d=1}^{L}b_{xd}}{\sum_{d=1}^{L}\max\{\ell_{xd},b_{xd}\}}.\]
For the finite-strata target, the broader stacked-program endpoints collapse to
\[H_{-,\mathcal X}=\frac{\sum_{x\in\mathcal X}\omega_x\sum_{d=1}^{L}(\ell_{xd}-a_{xd})_+}{\sum_{x\in\mathcal X}\omega_x\sum_{d=1}^{L}\min\{u_{xd},\max\{\ell_{xd},a_{xd}\}\}},\qquad
H_{+,\mathcal X}=\frac{\sum_{x\in\mathcal X}\omega_x\sum_{d=1}^{L}b_{xd}}{\sum_{x\in\mathcal X}\omega_x\sum_{d=1}^{L}\max\{\ell_{xd},b_{xd}\}}.\]
Both aggregate endpoints and every value between them are attainable by one covariate-augmented full-data law, and they require \(O(|\mathcal X|L)\) arithmetic operations.  The singleton-stratum case is \(\mathrm{thm{:}closed\mbox{-}form\mbox{-}harm\mbox{-}interval}\).  Thus the Cartesian prism and its ratio contraction are the closed-form collapse of the general KMV response-type program under KMV Strong Monotonicity.
%%% FIELD prop-ordered-collapse-proof
Fix a stratum \(x\).  Under \(\mathcal R_T^{\uparrow}\), the defining response constraints force \(p_{ij}=0\) for \(j<i\).  Since all flow coordinates are nonnegative, every expanded-state arc with a downward layer pair has zero mass, and every upper-triangular layer pair remains allowed.  Thus \(\Gamma_{\mathcal R_T^{\uparrow}}(\mu_x)\) is exactly the ordered expanded-state flow set in \(\mathrm{def{:}expanded\mbox{-}state\mbox{-}network}\), applied to the conditional observable law.  The decoder identifies
\[\mathsf S(\gamma_x)=\sum_{d=1}^{L}s_{xd},\qquad \mathsf N(\gamma_x)=\sum_{d=1}^{L}h_{xd}.\]
By \(\mathrm{thm{:}joint\mbox{-}stayer\mbox{-}harm\mbox{-}region}\), applied to that conditional law, the exact projection onto \((s_x,h_x)\) is the prism
\[\mathcal K_x=\prod_{d=1}^{L}\left\{(s_{xd},h_{xd}):\ell_{xd}\le s_{xd}\le u_{xd},\ (s_{xd}-a_{xd})_+\le h_{xd}\le\min\{s_{xd},b_{xd}\}\right\}.\]
The positive-stayer floor gives \(\sum_ds_{xd}\ge\kappa\).  The target theorem, applied conditionally, evaluates the minimum and maximum of \(\sum_dh_{xd}/\sum_ds_{xd}\) over \(\mathcal K_x\) as \(H_{-,x}\) and \(H_{+,x}\).  The general response-type theorem identifies the same extrema with the fixed-stratum flow program.  This proves the first display, including sharpness at zero cells.

Across finitely many strata, the decoder operates independently conditional on \(x\), so the exact feasible set of coordinates is \(\prod_x\mathcal K_x\).  For a fixed collection \((s_x)_x\), the smallest weighted ratio uses \(h_{xd}=(s_{xd}-a_{xd})_+\) and the largest uses \(h_{xd}=\min\{s_{xd},b_{xd}\}\).  Apply the coordinate-replacement proof of the target theorem to each coordinate \((x,d)\).  A replacement by \(\delta\) in stratum \(x\) changes the aggregate denominator, and when applicable its numerator, by \(\omega_x\delta\).  Thus the same inequalities apply with \(\delta\) replaced by \(\omega_x\delta\), and the optimizing coordinates are
\[s_{xd}^-=\min\{u_{xd},\max\{\ell_{xd},a_{xd}\}\},\qquad
s_{xd}^+=\max\{\ell_{xd},b_{xd}\}.\]
Substitution gives exactly \(H_{-,\mathcal X}\) and \(H_{+,\mathcal X}\).  Their denominators are at least
\[\sum_x\omega_x\kappa=\kappa>0.\]
The product prism is convex.  Mixing its two optimizing points with the denominator-adjusted coefficient used in the target proof attains every intermediate aggregate ratio.  The conditional decoder then combines the attaining stratum laws using the masses \(\omega_x\), producing one normalized global full-data law.  Finally, each of the \(|\mathcal X|L\) coordinates requires only a constant number of comparisons and arithmetic operations, proving the asserted complexity and the aggregate sharpness claim.
%%% FIELD target-proof
Fix \(P\in\mathcal D_{L,\kappa}\), and let
\[\Theta_I^H(P)=\left\{H(P^{\star}):P^{\star}\in\mathcal M_L\text{ induces the observable law }P\right\}\]
be the causal identified set.  This definition keeps the causal functional \(H(P^{\star})\) distinct from the observable endpoint maps to be derived.  Put
\[g_d(t)=(t-a_d)_+,qquad m_d(t)=\min\{t,b_d\}.\]
By \(\mathrm{thm{:}joint\mbox{-}stayer\mbox{-}harm\mbox{-}region}\), the exact image of all agreeing full-data laws under \(P^{\star}\mapsto(s,h)\) is
\[\mathcal K(P)=\left\{(s,h):\ell_d\le s_d\le u_d,\quad
g_d(s_d)\le h_d\le m_d(s_d),\quad d=1,\ldots,L\right\}.\]
This equality supplies both necessity and simultaneous attainability; it is stronger than coordinatewise outer bounds.  Moreover, for every \((s,h)\in\mathcal K(P)\),
\[S=\sum_{d=1}^{L}s_d\ge\sum_{d=1}^{L}\ell_d\ge\kappa>0\]
by \(\mathrm{ass{:}positive\mbox{-}stayer\mbox{-}floor}\).  Hence all ratios below are well defined, and
\[\Theta_I^H(P)=\left\{\frac{\sum_{d=1}^{L}h_d}{\sum_{d=1}^{L}s_d}:(s,h)\in\mathcal K(P)\right\}.\]

For a fixed \(s\), the denominator is positive and independent of \(h\), so the smallest ratio takes \(h_d=g_d(s_d)\) in every layer.  Define the feasible coordinate
\[s_d^- =\min\{u_d,\max\{\ell_d,a_d\}\}.\]
Starting from any \(s\in\prod_d[\ell_d,u_d]\), replace its coordinates one at a time by \(s_d^-\).  If \(s_d<s_d^-\), feasibility of \(s_d\) rules out \(s_d^-=\ell_d>s_d\), and the definition of the clipping point gives \(s_d\le s_d^-\le a_d\).  Thus the numerator contribution remains zero while the denominator weakly increases, so the ratio cannot increase.  If \(s_d>s_d^-\), the only possible clipping regimes give \(s_d\ge s_d^-\ge a_d\).  With \(\delta=s_d-s_d^-\), numerator and denominator both decrease by \(\delta\).  Writing their values before replacement as \(N\) and \(S\), the prism inequalities give \(0\le N\le S\), and
\[\frac{N-\delta}{S-\delta}\le\frac NS
\quad\Longleftrightarrow\quad
\delta(S-N)\ge0.\]
The new denominator remains at least \(\kappa\).  Therefore no coordinate replacement increases the ratio, and \(s^-\) is a global minimizer.  Directly from the three clipping regimes,
\[g_d(s_d^-)=(\ell_d-a_d)_+,\]
so
\[\inf\Theta_I^H(P)=
\frac{\sum_{d=1}^{L}(\ell_d-a_d)_+}
{\sum_{d=1}^{L}\min\{u_d,\max\{\ell_d,a_d\}\}}=H_-(\mu).\]

For the upper endpoint, the observable simplex restrictions in \(\mathrm{def{:}observable\mbox{-}mu\mbox{-}domain}\) imply
\[b_d\le q_{0d}\le r_d,qquad b_d\le c_d-q_{1d}\le c_d,\]
and hence \(0\le b_d\le u_d\).  For fixed \(s\), the largest ratio takes \(h_d=m_d(s_d)\).  Set
\[s_d^+=\max\{\ell_d,b_d\}\in[\ell_d,u_d].\]
If \(s_d<s_d^+\), feasibility rules out \(s_d^+=\ell_d>s_d\), so \(s_d^+=b_d\); increasing to \(b_d\) adds the same \(\delta\ge0\) to numerator and denominator.  Since \(N\le S\),
\[\frac{N+\delta}{S+\delta}\ge\frac NS
\quad\Longleftrightarrow\quad
\delta(S-N)\ge0.\]
If \(s_d>s_d^+\), both old and new coordinates are at least \(b_d\), so lowering to \(s_d^+\) leaves the numerator contribution equal to \(b_d\) and weakly decreases the positive denominator.  Thus \(s^+\) is a global maximizer.  Since \(m_d(s_d^+)=b_d\),
\[\sup\Theta_I^H(P)=
\frac{\sum_{d=1}^{L}b_d}
{\sum_{d=1}^{L}\max\{\ell_d,b_d\}}=H_+(\mu).\]

It remains to show that the identified set has no gaps.  The hinge description is equivalently the finite system of linear inequalities
\[\ell_d\le s_d\le u_d,qquad h_d\ge0,qquad h_d\ge s_d-a_d,qquad
h_d\le s_d,qquad h_d\le b_d,\]
so \(\mathcal K(P)\) is convex.  Let \((s^-,h^-)\) and \((s^+,h^+)\) be the two optimizing points above, and put \(S_-=\sum_ds_d^-\) and \(S_+=\sum_ds_d^+\), both at least \(\kappa\).  For any \(\lambda\in[0,1]\), define
\[t=\frac{\lambda S_-}{(1-\lambda)S_++\lambda S_-},\qquad
(s^{\lambda},h^{\lambda})=(1-t)(s^-,h^-)+t(s^+,h^+).\]
Convexity gives \((s^{\lambda},h^{\lambda})\in\mathcal K(P)\), and substitution yields
\[\frac{\sum_dh_d^{\lambda}}{\sum_ds_d^{\lambda}}
=(1-\lambda)H_-(\mu)+\lambda H_+(\mu).\]
The exact joint-region theorem lifts each such point, including both endpoints, to one normalized agreeing full-data law.  Consequently
\[\Theta_I^H(P)=[H_-(\mu),H_+(\mu)],\]
which proves sharpness.  Finally, each endpoint uses a constant number of additions, comparisons, and one final division per layer plus two sums.  Its arithmetic cost is therefore \(O(L)\).
